\section{Alcance, riesgos y criterios de salida}

Dentro del alcance: AMM constante, precio CEX externo, cuatro agentes, mapa discreto
local trazable, E0--E3, reproducibilidad, experimentos sintéticos y comparación limitada
con hechos de mercado.

Fuera del alcance: predicción rentable, exchanges operativos, réplicas LOB completas,
consenso blockchain, MEV avanzado, lending, bridges, aprendizaje por refuerzo y
garantías generales sobre todos los mercados. También queda fuera, por ahora, afirmar
validez empírica del mapa con base exclusiva en las simulaciones preliminares. Riesgo
principal: que el mapa reducido no aproxime adecuadamente al ABM fuera de una región
local. Mitigación: separar selección y confirmación, registrar configuraciones,
contrastar ablaciones y reportar desacuerdos como resultado. Otros riesgos:
sobreparametrización, dependencia de semillas, supuestos de liquidez ilimitada,
confusión entre estabilidad local y global, y sesgo en calibración.

\section{Discusión y límites de inferencia}

El análisis debe distinguir tres niveles: propiedades mecánicas deducidas del AMM ideal,
diagnóstico local del mapa reducido y patrones de un ABM discreto estocástico sujetos a
validación. La condición de Schur del mapa y la contracción preliminar no constituyen
garantías de estabilidad de mercados reales ni sustituyen una evaluación con datos
públicos. La investigación deberá reportar los regímenes donde la aproximación reduzca
bien el error y aquellos donde falle por saturación, límites de capital, choque
persistente o cambios de liquidez.

\section{Conclusiones y trabajo futuro}

El proyecto queda planteado como un marco experimental centrado en mecanismos:
seguimiento de precios, arbitraje limitado y fragilidad de liquidez. Las simulaciones
preliminares verifican que el prototipo reproduce las direcciones comparativas
esperadas para E0--E3 y permiten reemplazar una ODE genérica por un mapa discreto con
variables observables. Falta la fase decisiva: especificar criterios ex ante, ejecutar
sensibilidad y confirmación separadas, y evaluar la validez externa con datos públicos
acotados. Ninguno de los resultados actuales demuestra utilidad en el mercado real.
